% from ega0-1-fr.tex
\oldpage[0\textsubscript{I}]{13}
\begin{env}[1.2.1]
\label{0.1.2.1-fr}
On dit qu'une partie $S$ d'un anneau $A$ est \emph{multiplicative} si $1\in S$ et si le produit de deux éléments de $S$ est dans $S$.
Les exemples qui seront les plus importants pour la suite sont~: $1^\circ$ l'ensemble $S_f$ des puissances $f^n$ ($n\geq 0$) d'un élément $f\in A$; $2^\circ$ le complémentaire $A-\mathfrak{p}$ d'un idéal \emph{premier} $\mathfrak{p}$ de $A$.
\end{env}

\begin{env}[1.2.2]
\label{0.1.2.2-fr}
Soient $S$ une partie multiplicative d'un anneau $A$, $M$ un $A$-module; dans l'ensemble $M\times S$, la relation entre couples $(m_1,s_1)$, $(m_2,s_2)$~:
\[
  \text{<<~il existe $s\in S$ tel que $s(s_1m_2-s_2m_1)=0$~>>}
\]
est une relation d'équivalence.
On désigne par $S^{-1}M$ l'ensemble quotient de $M\times S$ par cette relation, par $m/s$ l'image canonique dans $S^{-1}M$ du couple $(m,s)$; on appelle \emph{application canonique} de $M$ dans $S^{-1}M$ l'application $i_M^S:m\mapsto m/1$ (aussi notée $i^S$).
Cette application n'est en général ni injective ni surjective; son noyau est l'ensemble des $m\in M$ tels qu'il existe un $s\in S$ pour lequel $sm=0$.

Dans $S^{-1}M$ on définit une loi de groupe additif en prenant
\[
  (m_1/s_1)+(m_2/s_2)=(s_2m_1+s_1m_2)/(s_1s_2)
\]
(on vérifie que c'est bien indépendant des expressions des éléments de $S^{-1}M$ considérés).
Sur $S^{-1}A$ on définit en outre une loi multiplicative en prenant $(a_1/s_1)(a_2/s_2)=(a_1a_2)/(s_1s_2)$, et enfin une loi externe sur $S^{-1}M$, ayant $S^{-1}A$ comme ensemble d'opérateurs, en posant $(a/s)(m/s')=(am)/(ss')$.
On vérifie ainsi que $S^{-1}A$ est muni d'une structure d'anneau (dit \emph{anneau de fractions de $A$ à dénominateurs dans $S$}) et $S^{-1}M$ d'une structure de $S^{-1}A$-module (dit \emph{module des fractions de $M$ à dénominateurs dans $S$}); pour tout $s\in S$, $s/1$ est inversible dans $S^{-1}A$, son inverse étant $1/s$.
L'application canonique $i_A^S$ (resp. $i_M^S$) est un homomorphisme d'anneaux (resp. un homomorphisme de $A$-modules, $S^{-1}M$ étant considéré comme $A$-module au moyen de l'homomorphisme $i_A^S:A\to S^{-1}A$).
\end{env}

\begin{env}[1.2.3]
\label{0.1.2.3-fr}
Si $S_f=\{f^n\}_{n\geq 0}$ pour un $f\in A$, on écrit $A_f$ et $M_f$ au lieu de $S_f^{-1}A$ et $S_f^{-1}M$; quand $A_f$ est considéré comme algèbre sur $A$, on peut écrire $A_f=A[1/f]$.
$A_f$ est isomorphe à l'algèbre quotient $A[T]/(fT-1)A[T]$.
Lorsque $f=1$, $A_f$ et $M_f$ s'identifient canoniquement à $A$ et $M$; si $f$ est nilpotent, $A_f$ et $M_f$ sont réduits à $0$.

Lorsque $S=A-\mathfrak{p}$, où $\mathfrak{p}$ est un idéal premier de $A$, on écrit $A_{\mathfrak{p}}$ et $M_{\mathfrak{p}}$ au lieu de $S^{-1}A$ et $S^{-1}M$; $A_{\mathfrak{p}}$ est un \emph{anneau local} dont l'idéal maximal $\mathfrak{q}$ est engendré par $i_A^S(\mathfrak{p})$, et on a $(i_A^S)^{-1}(\mathfrak{q})=\mathfrak{p}$; par passage aux quotients, $i_A^S$ donne un monomorphisme de l'anneau intègre $A/\mathfrak{p}$ dans le corps $A_{\mathfrak{p}}/\mathfrak{q}$, qui s'identifie au corps des fractions de $A/\mathfrak{p}$.
\end{env}

\begin{env}[1.2.4]
\label{0.1.2.4-fr}
L'anneau de fractions $S^{-1}A$ et l'homomorphisme canonique $i_A^S$ sont solution d'un \emph{problème d'application universelle}~: tout homomorphisme $u$ de $A$ dans un anneau $B$ tel que $u(S)$ se compose d'éléments \emph{inversibles} dans $B$ se factorise d'une seule manière
\[
  u:A\xrightarrow{i_A^S}S^{-1}A\xrightarrow{u^*}B
\]

